\lysh{22172_SOLUTION}{1}
\Lambda \Upsilon \Sigma \Eta

\alpha) Η απόσταση του σημείου $(x_o,y_o)$ από την ευθεία $Ax + By + \Gamma = 0$ δίνεται από τον τύπο:
\[ d = \frac{|Ax_o + By_o + \Gamma|}{\sqrt{A^2 + B^2}} \]
Επομένως έχουμε $d(A,\varepsilon) = \frac{|3 \cdot (-2) - 4 \cdot 1|}{\sqrt{3^2 + (-4)^2}} = \frac{|-10|}{\sqrt{25}} = \frac{10}{5} = 2$.

\beta) Η ζητούμενη ευθεία $(\eta)$ είναι κάθετη στην $(\varepsilon)$, οπότε το γινόμενο των συντελεστών διεύθυνσης της $(\eta)$ και της $(\varepsilon)$ θα είναι $-1$. Ο συντελεστής διεύθυνσης της $(\varepsilon)$ είναι
\[ \lambda_\varepsilon = -\frac{A}{B} = -\frac{3}{-4} = \frac{3}{4} \]
Επομένως ο συντελεστής διεύθυνσης της ζητούμενης ευθείας $(\eta)$ θα είναι $\lambda_\eta = -\frac{1}{\lambda_\varepsilon} = -\frac{1}{3/4} = -\frac{4}{3}$.

Η ευθεία διέρχεται από το $A(-2,1)$, οπότε η εξίσωση θα είναι $y-y_o = \lambda(x-x_o)$ ή
\[ y - 1 = -\frac{4}{3}(x+2) \]
ή
\[ y = -\frac{4}{3}x - \frac{8}{3} + 1 \]
ή
\[ y = -\frac{4}{3}x - \frac{5}{3} \]

\gamma) Η εξίσωση κύκλου με κέντρο το σημείο $(x_o,y_o)$ και ακτίνα $\rho$ δίνεται από την εξίσωση
\[ (x - x_o)^2 + (y - y_o)^2 = \rho^2 \quad (1) \]
Για να εφάπτεται ο κύκλος στην ευθεία $(\varepsilon)$, θα πρέπει η απόσταση του κέντρου του $A$ από την $(\varepsilon)$ να ισούται με την ακτίνα $\rho$. Στο ερώτημα ($\alpha$) βρήκαμε ότι η απόσταση του $A$ από την $(\varepsilon)$ είναι $2$, επομένως $\rho = 2$ και η (1) γίνεται
\[ (x - (-2))^2 + (y - 1)^2 = 2^2 \]
ή
\[ (x+2)^2 + (y-1)^2 = 4 \]
\end{lysh}